\begin{document}
\lhead[ \leftmark ]{Informatik, Freifach ``Applied Data Analysis''}
\rhead[Informatik]{Cyril Wendl \today, Winterthur}

\section*{Auftrag}

Im Rahmen dieses Freifachs erstellen Sie Ihre eigene interaktive Webseite über ein Thema nach Ihrer Wahl. Dabei lernen Sie verschieden Visualisierungsmöglichkeiten für grosse Datensets kennen. Wählen Sie zunächst ein spannendes Datenset und reichen Sie dieses (als Link) auf Moodle ein. Eine Auflistung nützlicher Links zu möglichen Datensets finden Sie auf Moodle

\subsection*{Anforderungen}

Ihre Webseite sollte dieses Thema auf schöne und interaktive Art zusammenfassen und visualisieren. Folgende Visualisierungen sollten auf sinnvolle Weise verwendet werden:
\begin{todolist}
    \item Ihre Webseite soll als \href{https://dash.plotly.com/tutorial}{\texttt{Dash}-Dashboard} umgesetzt werden. Alle Grafiken, Tabellen etc. sollen in die Webseite eingebettet sein.
    \item Der Code für Ihr Dashboard soll auf \href{https://github.com}{GitHub} gehostet werden.
    \item Ihre Dash-Webseite soll öffentlich auf \href{https://posit.co/blog/deploying-a-dash-application-to-posit-connect/}{Posit} gehostet sein.
    \item Ihr CSV-Datenset mit \href{https://pandas.pydata.org/docs/reference/api/pandas.read_csv.html}{\texttt{pandas}} eingelesen und verarbeitet werden. Achten Sie sich darauf, dass die Datei nicht zu gross ist (maximal einige 100 MB) und bereinigen Sie die Daten falls nötig mit einem separaten Python-Skript.
    \item Mindestens eine sinnvolle \href{https://pandas.pydata.org/docs/reference/api/pandas.DataFrame.html}{\texttt{DataFrame}}-Tabelle mit \href{https://dash.plotly.com/datatable}{\texttt{Dash DataTable}} anzeigen
    \item Mindestens vier sinnvolle \href{https://matplotlib.org/stable/tutorials/pyplot.html#sphx-glr-tutorials-pyplot-py}{\texttt{Matplotlib}}-Grafik und eine sinnvolle \href{https://seaborn.pydata.org/generated/seaborn.scatterplot.html}{\texttt{seaborn}}-Grafik anzeigen, z.B.:
    \begin{todolist}
        \item \texttt{matplotlib}: Z.B. \href{https://matplotlib.org/stable/tutorials/pyplot.html#sphx-glr-tutorials-pyplot-py}{Liniendiagramm (\textit{line plot})}, \href{https://matplotlib.org/stable/tutorials/pyplot.html#sphx-glr-tutorials-pyplot-py}{Streudiagramm (\textit{scatter plot})}, \footnote{Bei den Diagrammen sollten die Elemente wie Achsen, Titel, Legende, Gitterlinien, Farben usw. sinnvoll gewählt sein.}
        \item Sinnvolle Darstellung \textit{kategorischer Variablen}, z.B. mit \href{https://matplotlib.org/stable/tutorials/pyplot.html#sphx-glr-tutorials-pyplot-py}{Balkendiagramm (\textit{bar plot})}, \href{https://matplotlib.org/stable/tutorials/pyplot.html#sphx-glr-tutorials-pyplot-py}{Kreisdiagramm (\textit{pie chart})} oder \href{https://seaborn.pydata.org/generated/seaborn.boxplot.html}{Boxplot}
    \end{todolist}
    \item Ihre Webseite soll mindestens 5 interaktive Visualisierungen mit \href{https://plotly.com/python/}{plotly} enthalten.
    \item Ihre Webseite sollte inhaltlich sinnvoll aufgeteilt sein und für die unterschiedlichen thematischen Abschnitte sollten Sie \href{https://dash.plotly.com/dash-core-components/tabs}{Tabs}, \href{https://dash.plotly.com/dash-core-components/collapse}{Collapse-Komponenten} oder \href{https://dash.plotly.com/dash-core-components/dropdown}{Dropdown-Menüs} verwenden.
    \item Ihre \texttt{Dash}-Webseite sollte mindestens drei der folgenden Elemente sinnvoll einsetzen: \href{https://dash.plotly.com/dash-core-components/dropdown}{Dropdown-Menü}, \href{https://dash.plotly.com/dash-core-components/radioitems}{Radio-Buttons}, \href{https://dash.plotly.com/dash-core-components/checklist}{Checkboxen} oder \href{https://dash.plotly.com/dash-core-components/slider}{Slider}, mit welchen die Visualisierungen angepasst werden können.
    \item Ihre Webseite soll zusammenhänge innerhalb Ihres Datensets (Korrelationen) als \href{https://plotly.com/python/heatmaps/}{Korrelogramm} mit \texttt{plotly} darstellen
    \item Ihre Dash-Webseite sollte Angaben zum Autor (Ihnen) sowie zu der Quelle der Datensets enthalten.
\end{todolist}

\subsection*{Bonus}

\begin{todolist}
    \item Wenn Ihre Dash-Webseite \href{https://www.dash-bootstrap-components.com}{Bootstrap 5} verwendet.
    \item Wenn Ihre Dash-Webseite ein Kontaktformular enthält, mit dem NutzerInnen Nachrichten oder Feedback senden können.
\end{todolist}

\subsection*{Abgabe des Dashboards}
Wir werden am Schluss alle Webseiten auf eine öffentliche Web-Adresse auf \href{https://posit.co/blog/deploying-a-dash-application-to-posit-connect/}{Posit} hochladen. Reichen Sie dazu bitte ihre Posit-URL via \href{https://moodle.ksimlee.ch/}{Moodle} ein.

\subsection*{Prämierung der besten Webseiten}


\emoji{trophy} Die 2 Webseiten werden mit einem Preis belohnt! Alle TeilnehmerInnen, welche die obigen Anforderungen erfüllen, erhalten zudem ein Zertifikat am Ende des Kurses.
\end{document}